\section{Entity Relationship Model}

\subsection*{ER-Diagram Notation}

\emph{Entity} \hfill rectangle\\
\emph{Attribute} \hfill ellipse\\
\ind\emph{Primary Key} \hfill ellipse with underline\\
\ind\emph{Composite Attribute} \hfill ellipse with sub-ellipses\\
\ind\emph{Derived Attribute} \hfill dashed ellipse\\
\ind\emph{Multivalued Attribute} \hfill double ellipse\\
\emph{Relationship} \hfill diamond\\
\ind\emph{Role} \hfill text on line next to relationship\\
\hfill useful for unary relationships\\
\emph{Mapping Cardinality} \hfill\\
\ind\emph{One}\hfill arrow\\
\ind\emph{Many}\hfill no arrow\\
\emph{Participation} \hfill\\
\ind\emph{Total}\hfill double line\\
\ind\emph{Partial}\hfill single line\\
\emph{Weak Entity}\hfill\\
\ind\emph{Weak Entity Set}\hfill double rectangle\\
\ind\emph{Identifying Relationship}\hfill double diamond\\
\ind\emph{Discriminator}\hfill dashed underline\\
\emph{Specialization}\hfill\\
\ind\emph{Specialization}\hfill upside-down triangle with ``ISA''\\
\ind\emph{Disjoint}\hfill ``disjoint'' next to triangle\\
\ind\emph{Overlapping}\hfill -{}-\\
\ind\emph{Total Specialization}\hfill double line from ``parent''\\
\ind\emph{Partial Specialization}\hfill single line\\

\subsection*{Keys}

\emph{Super key}:\\\ind a set of 1+ attributes that uniquely\\\ind identifies a record\\
\emph{Candidate key}:\\\ind a minimal super key\\
\emph{Primary key}:\\\ind a candidate key chosen to identify records\\
\[
    \text{Primary Key} \subseteq \text{Candidate Key} \subseteq \text{Super Key}
\]
\emph{Foreign key}:\\\ind an attribute in a relation that is a primary\\\ind key in another relation\\

\subsection*{Weak Entity Relationship}

The weak entity set must have \emph{total participation} (ensures identifying entity exists), \emph{M-to-1/1-to-1} relationship with the identifying entity set (ensures unique identification).\\
{\hfill PK(Weak Entity)=PK(Identifying)+Discriminator \hfill}

\subsection*{ER-Diagrams to Schemas}
\emph{Regular entity \& attribute}: entity $\to$ table; attribute $\to$ column\\
\emph{Primary key}: underline attribute in schema\\
\emph{Composite attribute}: dot-access path to leaf component\\
\emph{Derived attribute}: not included in schema\\
\emph{Multivalued attribute}: separate entity, with FK to owning entity\\
\emph{Weak entity set}: FK to identifying entity set

\emph{Relationships}:\\
\ind\emph{M-to-M}: relationship $\to$ table; FKs to participating entities\\
\ind\emph{1-to-M/M-to-1}: FK in M-side entity to 1-side entity\\
\ind\emph{1-to-1}: FK in either entity to the other entity

\emph{Specialization}:\\
\ind\emph{(1)} higher entity $\to$ table; lower entity $\to$ table; PK of higher entity $\to$ FK in lower entity. \emph{[+ve]}: no redundancy; \emph{[-ve]}: less efficient retrieving\\
\ind\emph{(2)} lower entity table has all inherited attributes; PK of higher entity $\to$ FK in lower entity. \emph{[+ve]}: more efficient retrieving; \emph{[-ve]}: redundancy\\